\documentclass[conference]{IEEEtran}
\usepackage{amsmath,amssymb}
\usepackage{booktabs}
\usepackage{graphicx}
\usepackage[hidelinks]{hyperref}

\begin{document}

\title{Kernel Concentration and Bandwidth Scaling\\in Quantum Kernel Methods}

\author{\IEEEauthorblockN{Muhammad Ahmad Masood and Muhammad Waqas}
\IEEEauthorblockA{AI 631 Quantum Machine Learning, Fall 2026 --- Project \#5}}

\maketitle

\begin{abstract}
% DRAFT: written after the results are in.
\end{abstract}

\section{Introduction}
Quantum kernel methods map classical data into quantum states and use overlaps between states as a kernel for a classical learner, such as a support vector machine (SVM)~\cite{havlicek2019supervised,schuld2019feature}. Their appeal is that the feature space is an exponentially large Hilbert space. That same size is also a liability: as the number of qubits $n$ grows, the overlaps between states encoding different inputs can become exponentially small and nearly identical for every pair of inputs. When this happens, the kernel matrix approaches the identity and a model trained on it generalizes poorly~\cite{huang2021power}. Thanasilp \emph{et al.}\ called this effect \emph{exponential concentration}. They identified expressivity, entanglement, global measurements, and hardware noise as its sources, and they showed that a polynomial number of measurement shots is then not enough to tell kernel values apart~\cite{thanasilp2024concentration}.

A simple countermeasure is to scale the input features by a bandwidth $\lambda<1$ before encoding them. Shaydulin and Wild showed numerically that tuning this bandwidth can restore generalization for quantum kernels with many qubits~\cite{shaydulin2022bandwidth}, and Canatar \emph{et al.}\ analyzed how bandwidth controls generalization~\cite{canatar2022bandwidth}. Two practical questions remain. First, the kernel must be estimated from a finite number of shots on real hardware, and reducing $\lambda$ also moves every kernel value toward $1$, so there may be no bandwidth at which the kernel is both informative and measurable. Second, a strongly scaled-down quantum kernel may behave much like a classical kernel, so any benefit has to be measured against an equally well-tuned classical baseline.

\textbf{Research question.} How fast do quantum kernel values concentrate as the qubit count grows, and does a bandwidth hyperparameter restore learnability?

\textbf{Hypotheses.} \emph{H1:} at fixed $\lambda=1$, the off-diagonal kernel variance decreases exponentially in $n$. \emph{H2:} there is a decreasing bandwidth schedule $\lambda^{*}(n)$ under which the off-diagonal variance stays above the shot-noise level at a fixed shot budget, and test accuracy at $n=16$ stays statistically above that obtained at $\lambda=1$.

\textbf{Falsifier.} We reject H2 in our setting if, at $n=16$ and $N=10^{4}$ shots, every $\lambda$ in the search grid gives an off-diagonal spread that cannot be distinguished from the shot-noise spread (variance-ratio test, $\alpha=0.05$), or if the best tuned $\lambda$ does not give significantly higher test accuracy than $\lambda=1$ across fixed outer folds.

\textbf{Contributions.} We (i) measure off-diagonal kernel variance and test accuracy from 2 to 16 qubits over a bandwidth grid in exact simulation, (ii) determine where concentration becomes indistinguishable from finite-shot noise, (iii) repeat the variance measurements under a noise model derived from IBM device calibration data and on IBM quantum hardware at selected qubit counts, and (iv) compare every quantum result with an RBF-kernel SVM tuned with the same search protocol.

\section{Background}
% DRAFT: fidelity kernels; concentration results of [Thanasilp et al.]; bandwidth results of [Shaydulin & Wild, Canatar et al.]; finite-shot estimation.

\section{Methods}

\subsection{Data}
We use Fashion-MNIST~\cite{xiao2017fashion}, restricted to the binary pair T-shirt/top (label 0) vs.\ Pullover (label 2). For each of five seeds we draw a class-balanced subsample of 150 images per class without replacement from the official training split, giving 300 points per seed. To isolate the effect of $n$ and $\lambda$ from image structure, we also use a synthetic dataset: two isotropic Gaussian classes in $\mathbb{R}^{n}$ with means $\pm\frac{s}{2}\mathbf{u}$, where $\mathbf{u}$ is a random unit vector and $s=2$. Its Bayes-optimal accuracy is $\Phi(1)\approx0.841$ for every $n$.

\subsection{Folds and preprocessing}
The course's NB11 starter notebook was not available to us, so we implemented equivalent fixed folds, leakage-free preprocessing, baselines, ablations, and noise sweeps ourselves. For each seed, five stratified outer folds are generated once, and each outer training set is split into three stratified inner folds. The fold indices are stored on disk, every model reads the same files, and each result records a hash of its fold file. Preprocessing consists of standardization, PCA to $n$ components (Fashion-MNIST only), and a second standardization that gives every component unit variance. It is fit on the training rows of each inner or outer split and only then applied to the held-out rows. A unit test corrupts the held-out rows and checks that neither the fitted transforms nor the selected hyperparameters change.

\subsection{Quantum kernel}
Each input $\mathbf{x}\in\mathbb{R}^{n}$ is encoded on $n$ qubits as $\mathbf{z}=\lambda\mathbf{x}$ using an IQP-type $ZZ$ feature map~\cite{havlicek2019supervised} with $r=2$ layers and nearest-neighbour entanglement on a line:
\begin{align}
U(\mathbf{z}) &= \prod_{\ell=1}^{r} D(\mathbf{z})\,H^{\otimes n},\\
D(\mathbf{z}) &= \exp\Bigl(i\sum_{j=1}^{n} z_j Z_j + i\sum_{j=1}^{n-1} z_j z_{j+1} Z_j Z_{j+1}\Bigr).
\end{align}
The kernel is the state fidelity $k_\lambda(\mathbf{x},\mathbf{x}')=|\langle0^{n}|U^{\dagger}(\lambda\mathbf{x}')U(\lambda\mathbf{x})|0^{n}\rangle|^{2}$. Each layer has $n$ Hadamard gates, $n$ single-qubit $R_Z$ rotations, and $n-1$ two-qubit $R_{ZZ}$ rotations. The compute--uncompute circuit $U^{\dagger}(\lambda\mathbf{x}')U(\lambda\mathbf{x})$ therefore contains $4(n-1)$ $R_{ZZ}$ gates, for example 60 at $n=16$. Device-specific depths and two-qubit gate counts after transpilation are reported with the hardware results.

\textbf{Exact simulation.} Because $D$ is diagonal in the computational basis, we compute the statevectors exactly with a batched Walsh--Hadamard transform followed by an elementwise phase. Gram matrices are then computed from inner products of the statevectors. We validate this implementation against Qiskit statevector simulation of the same circuit; the two agree to within $10^{-10}$ for $n\in\{1,2,3,5\}$ and $\lambda\in\{0.05,0.7,2\}$.

\textbf{Concentration measure.} For a Gram matrix $K$ over $m$ points we report the mean and the unbiased sample variance of the $m(m-1)/2$ off-diagonal entries.

\textbf{Finite-shot estimation.} On an ideal device, each entry is estimated as the frequency of the all-zero outcome in $N$ shots of the compute--uncompute circuit, i.e.\ $\hat k=B/N$ with $B\sim\mathrm{Binomial}(N,k)$, so its variance is $k(1-k)/N$. Diagonal entries equal $1$ exactly. For $N\in\{10^{2},10^{3},10^{4},10^{5}\}$ we sample every entry of every training and test Gram matrix independently and repeat the full nested model selection on the sampled matrices, so that the bandwidth and $C$ are chosen under the same shot budget used for testing. We compare the true spread with the shot noise through the ratio $\rho_N=\sigma^{2}_{\lambda}/(\overline{k(1-k)}/N)$, where the bar averages over off-diagonal pairs. We test for any remaining spread with a binomial dispersion test. Under complete concentration, where every off-diagonal value equals a common $c$, the $M$ sampled off-diagonal entries are i.i.d.\ $\mathrm{Binomial}(N,c)/N$, so $D=(M-1)s^{2}/(\hat c(1-\hat c)/N)$ is approximately $\chi^{2}_{M-1}$. Here $s^{2}$ is their sample variance and $\hat c$ their mean. Rejecting at $\alpha=0.05$ means the measured spread cannot be explained by shot noise. The $\chi^{2}$ reference distribution is a large-sample approximation. It needs no assumption about how the true kernel values depend on each other, because under the null the only randomness is independent shot noise. In 2000 simulated null matrices per setting, the rejection rate at $\alpha=0.05$ was $0.048$ ($c=0.6$, $N=500$, $m=30$), $0.044$ ($c=0.99$, $N=10^{3}$, $m=60$), and $0.051$ ($c=0.95$, $N=10^{4}$, $m=60$).

\textbf{Falsifier.} At $n=16$ and $N=10^{4}$, a bandwidth counts as distinguishable from shot noise if the dispersion test rejects in more than half of the 25 outer folds (5 seeds $\times$ 5 folds). H2 is rejected if no bandwidth is distinguishable, or if the tuned-bandwidth test accuracy is not significantly higher than at $\lambda=1$ (one-sided Wilcoxon signed-rank test over the 25 paired folds, $\alpha=0.05$).

\textbf{Noise model.} Noisy simulations use a device model built from a saved calibration snapshot of an IBM backend: per-qubit $T_1$, $T_2$, and readout error, and per-gate error and duration for the native gates $\{\mathrm{CZ}, R_Z, \sqrt{X}, X\}$. Each gate applies thermal relaxation over its duration, followed by a depolarizing channel whose strength accounts for the rest of the reported gate error (for $\mathcal{E}(\rho)=(1-p)\rho+pI/d$, the average gate infidelity is $r=p(d-1)/d$). Measurements carry a symmetric readout error. A scale factor $s$ multiplies the depolarizing and readout error rates and divides $T_1$ and $T_2$. We use $s\in\{0,0.5,1,2\}$, where $s=0$ is an ideal device with shot noise only. On test circuits this model agrees with Qiskit Aer's built-in device model to within $0.02$ in kernel value. Circuits are transpiled to the backend's coupling map and native gates at optimization level 3 with a fixed seed. For each $n\in\{2,4,6,8,10\}$ and $\lambda\in\{2^{-4},2^{-2},2^{-1},1\}$, we estimate the kernel between $m=16$ Fashion-MNIST points (all 120 pairs plus the 16 diagonal circuits) with $4000$ shots per circuit. For a set of measured values we report the \emph{excess variance}: the sample variance minus the estimated shot-noise variance $\overline{\hat k(1-\hat k)}/(N-1)$, which is unbiased for the spread of the underlying noisy-device kernel. Its 95\% confidence interval comes from bootstrap resampling of points rather than pairs, because pairs that share a point are dependent.

\textbf{Hardware.} % DRAFT: fill in from results/hardware/<tag>/ after the run (backend, date, calibration, job ID, shots, QPU seconds).
The circuits for $n\in\{2,4,6\}$, $\lambda\in\{2^{-4},2^{-2},1\}$, and $m=10$ points (45 pairs and 10 diagonals per block) are transpiled against the live target of the chosen IBM backend and executed with the Qiskit Runtime Sampler. The same ISA circuits are then sampled with the same number of shots under the noise model built from the calibration in effect when the job ran. Device and model are therefore compared on identical circuits, qubits, and shot counts.

\subsection{Model selection and baselines}
All kernel models are SVMs with precomputed kernels. Hyperparameters are chosen by nested cross-validation: for each outer fold, the pair $(h, C)$, with $h$ the kernel bandwidth and $C\in\{0.1,1,10,100\}$, is selected by mean inner-fold validation accuracy. The model is then refit on the full outer training set and scored once on the outer test set. The quantum bandwidth is searched over $\lambda\in\{2^{-10},2^{-9},\dots,2^{1}\}$ (12 values, including the unscaled $\lambda=1$). The RBF baseline $k(\mathbf{x},\mathbf{x}')=\exp(-\gamma\|\mathbf{x}-\mathbf{x}'\|^{2})$ uses $\gamma=g/n$ with $g=\lambda^{2}$ over the same 12 values of $\lambda$. For small $\lambda$ the fidelity kernel behaves as $1-O(\lambda^{2}\|\mathbf{x}-\mathbf{x}'\|^{2})$, so this pairs the two grids point by point. Both searches use the same grid size and the same selection protocol. Because each component has unit variance, $g=1$ corresponds to scikit-learn's default \texttt{gamma="scale"}. A first run with narrower grids selected the lowest value in 16--36\% of folds, so the grids were extended downward before the main experiments. Further baselines are a linear-kernel SVM, $\ell_2$-regularized logistic regression with $C$ tuned on the same grid, and the majority-class rate. We report the mean and standard deviation of test accuracy over five outer folds and five seeds.

\section{Results}
% DRAFT

\section{Discussion}
% DRAFT

\section{Conclusion}
% DRAFT

\section*{Author Contributions}
% DRAFT

\bibliographystyle{IEEEtran}
\bibliography{refs}

\end{document}
